\label{sec:dimensions}

The four outcome dimensions of the O-series are chosen to span the major empirical
literature on redistricting and democratic representation. Table~\ref{tab:dimensions}
summarizes each dimension's definition, measurement unit, data source, and the
O-series paper that provides the primary analysis.

\begin{table}[ht]
\centering
\caption{Four Outcome Dimensions: O-Series Framework}
\label{tab:dimensions}
\begin{tabular}{lllll}
\toprule
Dimension & Outcome Variable & Unit & Data Source & Paper \\
\midrule
Competitiveness & Margin $< 10\%$ seats & Fraction of seats & VEST 2020 & O.1 \\
Turnout & Precinct VAP turnout & Percentage points & VEST 2020 & O.2 \\
Polarization & DW-NOMINATE gap & Score units & VoteView & O.3 \\
Const.-Rep.\ Distance & Mean driving time & Minutes & OSRM/OSM & O.4 \\
\bottomrule
\end{tabular}
\end{table}

\subsection{Dimension 1: Electoral Competitiveness}

A district is \textit{competitive} if the two-party vote margin in a given election
falls within 10 percentage points of the 50\% threshold. We use the two-party presidential
vote share as the competitiveness predictor rather than the congressional vote share, because
presidential vote is more stable across elections and less subject to incumbency effects.
The fraction of competitive seats in a state's congressional delegation is the aggregate
competitiveness measure.

Partisan gerrymandering systematically reduces competitiveness by design: the optimal
partisan map packs opposition voters into a small number of safe seats while cracking
remaining opposition voters into a larger number of seats where they lose by modest
margins. Both moves reduce competitive seats. Compactness disrupts this: compact
districts tend to include geographically proximate communities that are more
politically similar, creating more homogeneous districts that are either safe or
marginal---but the marginal districts are not chosen by the mapmaker.

The theoretical prediction for O.1 is that \textsc{bisect} plans produce more competitive
seats on average relative to enacted gerrymanders in states where the enacted map
has a documented partisan slant ($|\text{EG}| > 0.08$), but that the direction of the
competitiveness difference is unpredictable in states without documented partisan
gerrymandering. This asymmetric prediction is important: compact plans are not
maximally competitive. They are competitive-neutral, which means they produce
more competitive seats than gerrymandered plans but not necessarily more than the
maximum achievable under algorithmic competitive maximization.

\subsection{Dimension 2: Voter Turnout}

Voter turnout at the precinct level is a function of multiple factors: mobilization,
registration accessibility, partisan intensity, candidate quality, and---the channel
of interest here---district competitiveness. The ``stimulation hypothesis'' in electoral
participation research \citep{hays1997} posits that competitive elections increase
turnout by raising the perceived electoral stakes for individual voters. If redistricting
affects competitiveness (O.1), it may affect turnout through this channel.

The identification challenge for O.2 is separating the redistricting-induced competitiveness
variation from the confounding of partisan geography. A precinct moved from a safe district
into a competitive district may show increased turnout not because of the redistricting
change per se but because competitive districts attract more campaign resources and
candidate quality. We address this using a shift-share instrument: the precinct-level
predicted margin under the new map, constructed from the precinct's presidential vote
history and the new district boundaries, serves as the instrument for realized competitiveness.

\subsection{Dimension 3: Legislative Polarization}

Legislative polarization---the ideological distance between the parties' congressional
medians---has increased monotonically since the 1970s \citep{mccarty2006}. The
redistricting literature's claim is that partisan gerrymandering contributes to
polarization by reducing competitive seats, which in turn selects for more
ideologically extreme candidates in safe-seat primaries \citep{fiorina2008}. The
mechanism has two steps: (1) fewer competitive seats $\to$ more safe seats $\to$
more primary-determined general elections; (2) primary electorates are more ideologically
extreme than general electorates $\to$ extreme candidates win safe-seat primaries.

The empirical evidence for both steps is weaker than the theory suggests.
\citet{caugheyTausanovitch2017} find that redistricting affects state legislative
outcomes through ideological selection in some contexts but not universally.
\citet{mccarty2009} find that polarization has increased even in states without
gerrymandering, suggesting geographic sorting as an independent driver.
The O.3 analysis uses a DW-NOMINATE based approach with state-level controls
for geographic sorting to isolate the redistricting contribution.

\subsection{Dimension 4: Constituent-Representative Geographic Distance}

Geographic distance between constituents and their representative is a direct measure
of representational access that has received surprisingly little empirical attention.
The intuition is straightforward: a constituent living in a district drawn as a narrow
corridor connecting two distant communities faces a longer trip to the representative's
local district office than a constituent in a geographically compact district. This
matters because district office access is the primary channel through which constituents
raise constituent services issues (visa applications, Social Security disputes, veterans'
benefits, and similar casework).

We operationalize this as mean driving time from each census tract's population-weighted
centroid to the representative's primary district office address, as listed on house.gov.
OSRM (Open Source Routing Machine) with OpenStreetMap road network data provides the
routing computation. The O.4 analysis, the most novel in the O-series, documents a
correlation of $r = -0.67$ between district Polsby-Popper compactness and mean constituent
driving time across the 435 congressional districts, and shows that \textsc{bisect} plans
reduce mean driving time from 48 minutes (enacted gerrymanders) to 34 minutes.
